\chapter{Introdução}
A robótica tem ampliado significativamente sua presença na sociedade, com 
sistemas que interagem de forma cada vez mais próxima com os seres humanos. 
Desde drones de entrega e veículos autônomos até manipuladores industriais,
 a robótica se consolidou como uma área fundamental para o desenvolvimento 
 tecnológico. Embora cada categoria de robô possua particularidades operacionais,
  todas enfrentam desafios teóricos e práticos comuns, especialmente relativos
   à estimação de estado e ao controle.

O estado de um robô é definido por um conjunto de variáveis, como posição, 
velocidade, aceleração e atitude (orientação) que, quando integradas aos
 modelos cinemático e dinâmico do robô, descrevem seu movimento no espaço com
  precisão. Para obter essas grandezas, o uso de sensores é indispensável, 
  tanto sensores proprioceptivos quanto exteroceptivos

Para a condução dos testes experimentais deste trabalho, será adotado o robô 
quadrúpede Go2 (versão EDU), fabricado pela Unitree. Esta plataforma é equipada 
com diversos sensores, incluindo sensores de força de contato nas patas, LiDAR 
4D ultra-wide, câmera grande angular, unidade de medida inercial (IMU) e sensores 
de posição nos motores (encoders).

A locomoção de robôs com pernas exige um módulo de estimação de estados
 altamente confiável. Por essa razão, o presente projeto utiliza o estimador 
 MUSE (Multi-Sensor State Estimator). Esse sistema foi concebido como uma 
 evolução direta da arquitetura apresentada no trabalho "Proprioceptive Sensor 
 Fusion for Quadruped Robot State Estimation", entregando melhorias consideráveis
  na precisão e na robustez em relação ao seu antecessor.

Em relação ao middleware, embora o ROS 2 (Robot Operating System) seja o padrão 
predominante para o desenvolvimento de sistemas robóticos, este trabalho empregará
 o middleware DLS2 (Dynamic Legged Systems), criado pelo Instituto Italiano de
  Tecnologia (IIT). A escolha se justifica pelo fato de o DLS2 ser um ambiente 
  otimizado para a dinâmica de robôs quadrúpedes, o que o torna ideal para a atual
   aplicação.

Além da validação em hardware, o projeto abrangerá uma etapa em ambiente 
simulado. Os sensores do robô foram modelados no simulador Gazebo para replicar 
as condições reais de operação. A análise da telemetria e dos dados do sistema 
será feita com o auxílio do software PlotJuggler, enquanto o processamento das 
nuvens de pontos do LiDAR e o mapeamento ficarão a cargo da biblioteca PCL 
(Point Cloud Library).
